% ============================================================================
% Fichier  : partie3/P03_C13_deontologie_investigateur.tex
% Titre    : Déontologie de l'Investigateur Numérique
% Auteur   : MINKA MI NGUIDJOÏ Thierry Emmanuel
% Version  : 1.0 -- Juin 2026
% Statut   : RÉDIGÉ
% Sources  : Compétences_non_techniques_IN_101250.pdf (29 pages -- CNT-7
%            droits, CNT-8 confidentialité, CNT-9 éthique, fondements
%            théoriques Kant/Mill/Aristote, biais cognitifs)
%            M3_Cadre_Normatif_National_PNC.docx (responsabilité OPJ)
%            Ch. PIU (ACPO Principe 4), Ch. custody, Ch. droit camerounais
% Positionnement : Ce chapitre clôt la Partie III. Il traite la dimension
%                  humaine du métier -- ce que les textes de loi ne couvrent
%                  pas. La déontologie n'est pas un supplément d'âme : c'est
%                  la condition de la confiance institutionnelle sur laquelle
%                  repose toute la chaîne judiciaire.
% ============================================================================

\chapter{Déontologie de l'Investigateur Numérique}
\label{chap:deonto}

\epigraph{%
    « L'excellence investigative n'est pas seulement maîtrise de techniques,
    mais développement d'un caractère vertueux~: intègre, courageux,
    prudent, juste. »%
}{--- \textit{Compétences Non Techniques de l'Investigateur Numérique},
LIMSI, 2025}

\niveauChapitre{Fondamental}{%
    Aucun prérequis technique. Ce chapitre est philosophique, éthique
    et pratique. Il complète les chapitres juridiques en abordant ce
    que les textes de loi ne couvrent pas~: les choix moraux que
    l'investigateur doit faire sous pression.%
}

\begin{boxobjectifs}
\begin{itemize}
    \item Comprendre pourquoi la déontologie est une condition de
          l'efficacité investigative, pas seulement une contrainte éthique.
    \item Maîtriser les trois cadres éthiques fondamentaux~: déontologique
          (Kant), conséquentialiste (Mill) et des vertus (Aristote).
    \item Identifier les cinq tensions déontologiques majeures auxquelles
          tout investigateur sera confronté~; savoir les résoudre.
    \item Reconnaître les biais cognitifs qui compromettent l'objectivité
          et appliquer des contre-mesures systématiques.
    \item Maîtriser les obligations de confidentialité et leurs limites
          légales en droit camerounais.
    \item Conduire un entretien investigatif respectueux des droits
          fondamentaux et juridiquement recevable.
\end{itemize}
\end{boxobjectifs}

\bigskip

% ============================================================================
\section*{Paradoxe de l'investigation numérique}
% ============================================================================

Les statistiques professionnelles sont instructives~:
\begin{itemize}
    \item 68\% des investigations numériques échouent non par défaillance
          technique mais par mauvaise gestion de la dimension humaine
          (SANS Institute, 2024)~;
    \item 73\% des preuves numériques rejetées en justice le sont pour
          vices procéduraux liés aux entretiens~: non-respect des droits,
          contamination testimoniale~;
    \item 82\% des investigateurs seniors identifient les \textit{soft skills}
          comme facteur différenciant majeur entre investigateurs compétents
          et investigateurs excellents.
\end{itemize}

La déontologie n'est donc pas un supplément d'âme~: c'est la condition
de la \textbf{confiance institutionnelle} sur laquelle repose toute
la chaîne judiciaire, et l'un des facteurs les plus prédictifs du
succès à long terme d'une carrière d'investigateur.

\separateur

% ============================================================================
\section{Fondements Éthiques~: Trois Cadres pour Décider}
\label{sec:fondements-ethiques}
% ============================================================================

\subsection{Éthique déontologique~: Kant et les droits absolus}

L'\termecle{éthique déontologique}\index{éthique déontologique} (Kant,
\textit{Fondements de la Métaphysique des M\oe urs}, 1785) évalue une
action selon sa conformité à un devoir ou une règle, indépendamment
des conséquences.

\begin{boxdef}
\textbf{Impératif catégorique de Kant}~: « Agis uniquement d'après
la maxime qui fait que tu peux vouloir en même temps qu'elle devienne
une loi universelle. »

\textbf{Droits inaliénables kantiens~:} la dignité humaine est non
négociable. L'humain est fin en soi, jamais simple moyen. Interdit
l'instrumentalisation~: manipuler une personne comme un outil viole
sa dignité, même si c'est efficace.

\textbf{Application investigative~:} un suspect a des droits indépendamment
de sa culpabilité présumée. L'interrogatoire coercitif qui viole ces
droits est éthiquement interdit kantiennement~--- même si les aveux
obtenu sont vrais.
\end{boxdef}

\subsection{Éthique conséquentialiste~: Mill et le calcul des conséquences}

L'\termecle{éthique conséquentialiste}\index{utilitarisme} (Mill,
\textit{L'Utilitarisme}, 1863) évalue la moralité d'une action par
ses conséquences~: maximiser le bien-être, minimiser la souffrance.

\begin{boiteRemarque}{Tension fondamentale Kant/Mill en investigation}
\textbf{Kant~:} droits absolus, non négociables (déontologie).\\
\textbf{Mill~:} droits relatifs au bien-être global (conséquences).

\textbf{Cas concret~:} un suspect refuse de déverrouiller son téléphone.
Des vies sont peut-être en danger.
\begin{itemize}
    \item Kant~: le droit à ne pas s'auto-incriminer est absolu.
          La coercition est interdite.
    \item Mill~: si les bénéfices (vies sauvées) l'emportent sur
          les coûts (violation d'un droit)~, la coercition pourrait
          être justifiée.
\end{itemize}
\textbf{Réponse légale camerounaise~:} art.~88 Loi~2010/012 tranche
la tension~: le refus après réquisition judiciaire est une infraction
autonome. La loi choisit la voie kantienne~--- la contrainte légale
remplace la coercition directe.
\end{boiteRemarque}

\subsection{Éthique des vertus~: Aristote et le caractère de l'investigateur}

L'\termecle{éthique des vertus}\index{éthique des vertus} (Aristote,
\textit{Éthique à Nicomaque}) demande non pas ``quelle règle dois-je
suivre~?'' mais ``quel type de personne dois-je être~?''

Les \textbf{quatre vertus de l'investigateur}~:
\begin{enumerate}
    \item \textbf{Intégrité}~: dire la vérité même quand elle contredit
          ses hypothèses initiales~; documenter les preuves qui réfutent
          autant que celles qui confirment.
    \item \textbf{Courage}~: signaler une violation de procédure commise
          par un collègue~; résister aux pressions hiérarchiques pour
          ``conclure vite''~; accepter l'incertitude.
    \item \textbf{Prudence} (\textit{phronèsis})~: la sagesse pratique
          qui permet de juger correctement dans les situations singulières.
          Aucun manuel n'anticipe toutes les situations~; la prudence
          est ce qui comble le vide.
    \item \textbf{Justice}~: traiter chaque personne équitablement,
          indépendamment de ses antécédents, de son appartenance sociale
          ou de la présomption de culpabilité.
\end{enumerate}

La trajectoire de développement déontologique comporte trois stades~:
\begin{description}
    \item[Stade~1 --- Conformité externe] Le débutant suit les règles
          mécaniquement (checklists, protocoles).
    \item[Stade~2 --- Intériorisation] L'intermédiaire comprend les
          raisons des règles et les applique avec jugement.
    \item[Stade~3 --- Vertu intégrée] L'expert a internalisé les vertus~;
          il agit vertueusement spontanément (``seconde nature'').
\end{description}

% ============================================================================
\section{Les Cinq Tensions Déontologiques Majeures}
\label{sec:tensions}
% ============================================================================

\subsection{Tension~1~: Efficacité vs Droits fondamentaux}

La pression de ``résoudre vite'' une affaire peut pousser à contourner
les procédures. Chaque raccourci procédural~:
\begin{enumerate}
    \item Viole les droits de la personne visée~;
    \item Fragilise la valeur probatoire des preuves obtenues~;
    \item Expose l'investigateur à des poursuites personnelles.
\end{enumerate}

\begin{boxCRO}
\textbf{Analyse CRO --- Tension efficacité vs droits}

$\mathcal{R}$~: les méthodes coercitives produisent des aveux non fiables
(taux de faux aveux~: 15--25\% selon les études). L'efficacité à court
terme se paye d'une fiabilité dégradée.

$\mathcal{O}$~: toute preuve obtenue en violation des droits fondamentaux
sera contestée et potentiellement exclue. L'opposabilité est maximisée
par le respect strict des droits, pas par leur contournement.

$\mathcal{C}$~: les méthodes coercitives violent la confidentialité
des données personnelles du suspect.

\cro{$\downarrow$ coercition viola privauté}{$\downarrow\downarrow$ aveux non fiables}{$\downarrow\downarrow$ exclusion probable}

\textbf{Conclusion~:} le respect des droits et l'efficacité ne sont
pas en tension --- ils convergent. Un dossier propre est un dossier solide.
\end{boxCRO}

\subsection{Tension~2~: Confidentialité vs Obligation de divulgation}

L'investigateur est soumis à deux obligations contradictoires~:
protéger les informations confidentielles qu'il collecte, ET divulguer
ce qui est nécessaire à la justice.

\begin{tableaucours}{p{3.8cm}p{4.0cm}p{4.0cm}}{%
    \tableauHeader{Situation} &
    \tableauHeader{Obligation} &
    \tableauHeader{Base légale}
}
Données personnelles collectées
    pendant l'investigation &
    Confidentialité stricte~; usage limité à la finalité de l'enquête &
    Art.~74 Loi~2010~; Art.~9--11 Loi~2024 \\
Preuves pertinentes pour la défense
    (innocence du suspect) &
    Obligation de divulgation au parquet~; ne pas dissimuler &
    CPP camerounais~; ACPO Principe~4 \\
Informations compromettantes
    sur des collègues &
    Obligation déontologique de signalement~; résistance à la
    complaisance corporative &
    Art.~89 Code pénal (recel)~; éthique professionnelle \\
Données protégées
    par le secret professionnel &
    Saisie limitée~; respect du secret médical, de l'avocat &
    CPP~; art.~5 Loi~2024 (données sensibles) \\
\end{tableaucours}
\vspace{-0.8em}
\begin{center}{\small\itshape Tableau~13.1 --- Tensions confidentialité / divulgation}\end{center}

\subsection{Tension~3~: Objectivité vs Biais de confirmation}

Le \termecle{biais de confirmation}\index{biais de confirmation}
--- la tendance à privilégier les informations qui confirment nos
croyances initiales --- est l'ennemi structurel de l'objectivité
investigative. Sa particularité~: il opère de façon \emph{inconsciente}.

\begin{boiteRemarque}{Le biais de confirmation en chiffres}
Une hypothèse préalable (``c'est lui le coupable'') modifie
l'interprétation de tout ce qui suit~:
\begin{itemize}
    \item Un interviewé nerveux~$\to$ s'il est étiqueté ``menteur''~:
          confirmation du mensonge. S'il est étiqueté ``victime''~:
          compréhension du traumatisme.
    \item Une incohérence dans le récit~$\to$ si hypothèse coupable~:
          preuve de mensonge. Si hypothèse innocent~: erreur de mémoire.
    \item Chaque comportement réinterprété à travers le filtre de
          l'étiquette initiale~--- biais de confirmation en cascade.
\end{itemize}
\textbf{Contre-mesure~:} E9bis du PIU (falsification active). Pour
chaque hypothèse~: chercher activement ce qui la réfuterait, pas
ce qui la confirme.
\end{boiteRemarque}

\subsection{Tension~4~: Loyauté institutionnelle vs Intégrité personnelle}

Quand un supérieur hiérarchique demande de ``conclure vite''~; quand
un collègue a contaminé une scène~; quand des preuves pointent vers
des personnes influentes~:

\begin{itemize}
    \item \textbf{Signalement interne}~: documenter précisément les
          faits~; remonter par la chaîne hiérarchique~; conserver une
          copie personnelle de toute documentation transmise.
    \item \textbf{Signalement externe}~: si le signalement interne
          échoue ou est bloqué~: ANTIC~; parquet~; inspection générale
          des services.
    \item \textbf{Protection~:} un OPJ qui signale de bonne foi une
          irrégularité bénéficie en principe d'une protection contre
          les représailles~--- même si la protection légale explicite
          est encore faible en droit camerounais.
\end{itemize}

\begin{boxalert}
\textbf{Le courage déontologique~: la compétence la moins enseignée}

Il est plus difficile de refuser d'exécuter un ordre illégal d'un
supérieur que de résister à un suspect hostile. Les études sur les
erreurs judiciaires montrent que la plupart impliquent au moins un
acteur qui ``savait'' mais s'est tu. Le courage déontologique est
la vertu qui protège la chaîne judiciaire de la corruption systémique.
\end{boxalert}

\subsection{Tension~5~: Compétence personnelle vs Délégation nécessaire}

L'investigateur doit savoir reconnaître les limites de ses propres
compétences. Intervenir au-delà de sa formation~:
\begin{itemize}
    \item Contamine les preuves (violation d'ACPO Principe~2)~;
    \item Engage sa responsabilité personnelle~;
    \item Fragilise l'ensemble du dossier.
\end{itemize}

\textbf{Règle pratique~:} « Si je ne sais pas faire sans risque de
modifier les données, je ne fais pas. J'appelle un spécialiste et
je documente pourquoi. » Cette auto-évaluation honnête est elle-même
une compétence déontologique.

% ============================================================================
\section{Droits des Personnes Interrogées}
\label{sec:droits-personnes}
% ============================================================================

\subsection{Cadre légal et principe}

Tout entretien investigatif implique des personnes~: victimes, témoins,
suspects. Chaque catégorie dispose de droits spécifiques qui conditionnent
la valeur probatoire des témoignages recueillis.

\begin{boxjuris}
\textbf{Droits fondamentaux en contexte d'entretien investigatif~:}
\begin{itemize}
    \item \textbf{Droit à l'information}~: toute personne entendue
          doit être informée de l'objet de l'entretien et de son
          droit de refuser (sauf obligation légale contraire).
    \item \textbf{Droit à l'assistance}~: un suspect en garde à vue
          a le droit à l'assistance d'un avocat dès le début
          (CPP camerounais).
    \item \textbf{Droit au silence}~: le suspect ne peut être contraint
          à s'auto-incriminer (principe \textit{nemo tenetur}).
    \item \textbf{Droits des mineurs}~: présence obligatoire d'un
          représentant légal~; protocoles spécifiques~; attention
          particulière aux processus de mémoire (contamination facile).
    \item \textbf{Droit à l'interprète}~: si la personne ne maîtrise
          pas la langue de l'entretien~; l'interprète doit être
          assermenté pour les actes judiciaires.
\end{itemize}
\end{boxjuris}

\subsection{La méthode PEACE~: entretien éthique et efficace}

La méthode \termecle{PEACE}\index{méthode PEACE}
(\textit{Preparation, Engage, Account, Closure, Evaluation}), développée
par la police britannique en 1992, est la référence mondiale pour les
entretiens investigatifs respectueux des droits~:

\begin{tableaucours}{p{2.8cm}p{8.8cm}}{%
    \tableauHeader{Phase} &
    \tableauHeader{Contenu et objectifs}
}
\textbf{P --- Préparation} &
    Définir l'objectif de l'entretien. Recueillir toutes les informations
    disponibles. Choisir le lieu et le moment adaptés. Préparer les
    questions ouvertes. \\
\textbf{E --- Engage and Explain} &
    Établir le rapport~; présenter l'objet~; informer des droits~;
    mettre à l'aise. Un témoin à l'aise révèle davantage. \\
\textbf{A --- Account} &
    Recueillir le récit libre sans interruption~; puis poser des
    questions de clarification. Ne jamais diriger~; ne pas suggérer. \\
\textbf{C --- Closure} &
    Résumer~; donner à l'interviewé l'opportunité de corriger,
    compléter ou ajouter. Informer des suites. \\
\textbf{E --- Evaluation} &
    Documenter immédiatement~; évaluer la cohérence~; identifier
    les lacunes~; planifier les entretiens suivants. \\
\end{tableaucours}
\vspace{-0.8em}
\begin{center}{\small\itshape Tableau~13.2 --- Méthode PEACE~: entretien éthique et efficace}\end{center}

\begin{boiteRemarque}{Pourquoi PEACE est plus efficace que l'interrogatoire coercitif}
Les études de psychologie légale montrent que les techniques d'entretien
basées sur l'instauration d'une confiance et la narration libre obtiennent
\emph{plus} d'informations utiles que les techniques coercitives, tout en
produisant \emph{moins} de faux aveux. L'entretien éthique n'est pas un
sacrifice d'efficacité~--- c'est une méthode supérieure.
\end{boiteRemarque}

% ============================================================================
\section{Biais Cognitifs et Contre-Mesures}
\label{sec:biais}
% ============================================================================

\subsection{Les six biais les plus dangereux en investigation}

\begin{tableaucours}{p{3.5cm}p{4.0cm}p{4.5cm}}{%
    \tableauHeader{Biais} &
    \tableauHeader{Manifestation investigative} &
    \tableauHeader{Contre-mesure}
}
\textbf{Confirmation} &
    Chercher des preuves qui confirment l'hypothèse initiale~;
    ignorer les preuves contraires &
    E9bis PIU~: chercher activement ce qui réfuterait l'hypothèse \\
\textbf{Ancrage} &
    Première information reçue oriente toutes les suivantes~;
    première suspicion difficile à réviser &
    Générer plusieurs hypothèses (E8 PIU) avant de fixer une conclusion \\
\textbf{Étiquetage prématuré} &
    Qualifier ``suspect'' ou ``menteur'' oriente l'interprétation
    de tous les comportements ultérieurs (Becker, 1963) &
    Maintenir des catégories ouvertes~; consigner les observations
    avant les interprétations \\
\textbf{Disponibilité} &
    Surestimer la probabilité d'un scénario parce qu'il est familier~;
    ``ce ressemble à ce que j'ai vu l'an dernier'' &
    Vérifier la ressemblance par les faits, pas par l'intuition~;
    raisonner sur les statistiques \\
\textbf{Surconfiance} &
    Croire être objectif quand on ne l'est pas~; ``je sais
    quand quelqu'un ment'' &
    Taux de précision humain dans détection du mensonge~:
    54\%~--- à peine mieux que le hasard \\
\textbf{Tunnelisation} &
    Se focaliser sur un seul suspect/scénario~; négliger
    les alternatives &
    ``Quelle est l'hypothèse que je n'ai pas explorée~?''
    Rotation de l'hypothèse active tous les 48h \\
\end{tableaucours}
\vspace{-0.8em}
\begin{center}{\small\itshape Tableau~13.3 --- Biais cognitifs et contre-mesures}\end{center}

\subsection{La pratique de l'objectivité~: protocoles concrets}

\begin{enumerate}
    \item \textbf{Séparer constatations et interprétations~:}
          dans chaque document~: une section ``Constatations''
          (faits vérifiables) et une section ``Analyse''
          (interprétations, avec niveau de certitude explicite).
          Jamais mêlés.
    \item \textbf{Documenter les preuves contraires~:} noter explicitement
          chaque preuve qui ne s'accorde pas avec l'hypothèse principale.
          Si le dossier n'en contient aucune, chercher plus.
    \item \textbf{Rotation des hypothèses~:} tous les deux jours, s'obliger
          à défendre l'hypothèse inverse. Que faudrait-il pour qu'elle
          soit vraie~?
    \item \textbf{Revue par les pairs~:} présenter les conclusions
          préliminaires à un collègue avec le mandat explicite de
          les contester.
    \item \textbf{Journal de bord éthique~:} noter les situations
          où la pression a été ressentie pour orienter les conclusions~;
          les moments de doute~; les décisions difficiles. Ce journal
          constitue une protection en cas de contestation ultérieure.
\end{enumerate}

% ============================================================================
\section{Responsabilité Personnelle de l'Investigateur}
\label{sec:responsabilite}
% ============================================================================

\subsection{Responsabilité pénale personnelle}

L'investigateur numérique n'est pas protégé par ``j'exécutais un ordre''.
En droit camerounais~:

\begin{itemize}
    \item \textbf{Responsabilité pour violation de procédure~:}
          un OPJ qui conduit une perquisition sans autorisation judiciaire
          (art.~52 Loi~2010) commet une faute personnelle~; il peut être
          poursuivi pour violation de domicile, atteinte à la vie privée.
    \item \textbf{Responsabilité pour faux~:} un investigateur qui
          falsifie un rapport d'expertise ou un formulaire de custody
          commet le crime de faux en écriture (art.~155--162 Code pénal)~;
          c'est une infraction personnelle, indépendante du résultat
          de l'affaire.
    \item \textbf{Responsabilité pour complicité~:} un investigateur
          qui couvre la violation procédurale d'un collègue peut être
          poursuivi pour complicité ou recel (art.~89 Code pénal).
\end{itemize}

\subsection{ACPO Principe~4~: responsabilité du chef d'investigation}

Le Principe~4 du guide ACPO établit que \textbf{le responsable de
l'investigation est personnellement responsable de la conformité de
l'ensemble de l'équipe}. Cette responsabilité~:
\begin{itemize}
    \item S'étend à toutes les actions de chaque membre de l'équipe~;
    \item Implique une supervision active~: pas de délégation aveugle~;
    \item Exige une documentation systématique des décisions~: chaque
          choix procédural doit pouvoir être justifié.
\end{itemize}

\subsection{L'obligation d'auto-évaluation continue}

Le formulaire d'auto-évaluation de compétences non techniques (grille
CNT~1--15) est un outil de développement professionnel~: identifier
ses lacunes avant qu'elles ne causent une erreur judiciaire. Les cinq
questions clés à se poser à la clôture de chaque investigation~:

\begin{enumerate}
    \item Ai-je cherché activement des preuves contraires à mon hypothèse~?
    \item Ai-je respecté scrupuleusement les droits de chaque personne
          rencontrée~?
    \item Ai-je documenté tous les transferts d'evidence, même temporaires~?
    \item Y a-t-il un moment où j'ai ressenti de la pression pour
          orienter mes conclusions~? L'ai-je documenté~?
    \item Mon rapport est-il défendable mot à mot en contre-interrogatoire~?
\end{enumerate}

% ============================================================================
\section{Déontologie et Témoignage Expert}
\label{sec:temoignage-expert}
% ============================================================================

L'investigateur peut être amené à témoigner en tant qu'expert devant
un tribunal. Cette situation exige une préparation spécifique.

\subsection{Obligations de l'expert judiciaire}

\begin{itemize}
    \item \textbf{Impartialité}~: l'expert n'est ni l'avocat de
          l'accusation ni celui de la défense. Il est le serviteur
          de la vérité~--- même si ses conclusions déplaisent à
          l'autorité qui l'a mandaté.
    \item \textbf{Compétence déclarée}~: l'expert ne doit témoigner
          que dans le domaine de sa compétence réelle. Affirmer une
          certitude que la technique ne permet pas est une faute grave.
    \item \textbf{Limites explicites}~: déclarer ce que l'analyse
          \emph{ne peut pas} établir est aussi important que déclarer
          ce qu'elle établit.
    \item \textbf{Distinction constatation/interprétation}~: devant
          le tribunal~: « J'ai constaté X. Une interprétation de X
          serait Y, mais d'autres explications sont possibles. »
\end{itemize}

\subsection{Résistance au contre-interrogatoire}

\begin{boiteRemarque}{Comment se préparer au contre-interrogatoire}
\begin{enumerate}
    \item Relire l'intégralité du rapport d'expertise la veille~;
    \item Identifier soi-même les trois points les plus attaquables~;
    \item Préparer une réponse factuelle et documentée pour chacun~;
    \item Pendant l'interrogatoire~: répondre aux questions posées,
          pas aux questions redoutées~; ne pas sur-interpréter~;
          dire ``je ne sais pas'' quand c'est vrai.
    \item Ne jamais modifier verbalement une conclusion écrite
          sous la pression d'un avocat compétent~: ``Mon rapport
          parle pour lui-même.''
\end{enumerate}
\end{boiteRemarque}

% ============================================================================
\section{Cas Pratique~: L'Affaire ONAMBELE --- Décisions Déontologiques}
\label{sec:deon-onambele}
% ============================================================================

\begin{boxscenario}{Affaire ONAMBELE --- trois décisions sous pression}
\textbf{Contexte~:} L'OPJ NKENG, en charge de l'affaire ONAMBELE (fraude
sur les marchés publics, Chapitre~\ref{chap:gr-affaires}), est confronté
à trois situations de pression déontologique au cours de la même semaine.

\textbf{Situation~1 --- Pression hiérarchique~:}
Le directeur de la DCPJ convoque l'OPJ NKENG et lui demande de ``clore
le dossier rapidement''~--- un responsable politique est dans l'entourage
du suspect. NKENG doit décider~: céder à la pression ou maintenir son
investigation.

\textbf{Analyse déontologique~:}
\begin{itemize}
    \item Vertu en jeu~: \textbf{courage}~--- la vertu la moins
          enseignée mais la plus nécessaire.
    \item Action correcte~: documenter précisément la demande (date,
          heure, contenu, nom du demandeur)~; poursuivre l'investigation
          selon les règles~; escalader au parquet si la pression s'intensifie.
    \item Risque~: représailles. Protection~: la documentation.
\end{itemize}

\textbf{Situation~2 --- Preuve contraire à l'hypothèse~:}
L'analyse forensique révèle que les virements suspects ont été initiés
depuis un compte que l'OPJ NKENG avait classé comme ``témoin innocent''~---
pas depuis le compte du suspect principal. L'hypothèse principale s'effondre.

\textbf{Analyse déontologique~:}
\begin{itemize}
    \item Biais en jeu~: \textbf{biais de confirmation}~--- résistance
          à réviser une hypothèse investie de plusieurs semaines de travail.
    \item Action correcte~: consigner immédiatement la découverte dans
          le journal de bord~; rouvrir le champ des hypothèses (E8--E9
          PIU)~; informer le parquet de la révision d'hypothèse.
    \item Vertu en jeu~: \textbf{intégrité}~--- dire la vérité même
          quand elle contrarie son travail antérieur.
\end{itemize}

\textbf{Situation~3 --- Violation par un collègue~:}
L'OPJ NKENG découvre que son collègue BITA a interrogé un suspect sans
présence d'avocat et sans l'informer de ses droits~--- les déclarations
obtenues sont potentiellement irrecevables.

\textbf{Analyse déontologique~:}
\begin{itemize}
    \item Tension~: loyauté envers un collègue vs intégrité du dossier.
    \item Action correcte~: informer le parquet~; documenter la violation~;
          proposer de faire réentendre le suspect dans le respect de
          ses droits.
    \item Conséquence de l'inaction~: si la violation n'est pas signalée
          et est découverte par la défense à l'audience, c'est
          l'ensemble du dossier qui est fragilisé --- et l'OPJ NKENG
          qui a couvert peut être impliqué.
\end{itemize}
\end{boxscenario}

\separateur

\begin{boxresume}
\textbf{Ce qu'il faut retenir de ce chapitre~:}
\begin{itemize}
    \item La déontologie n'est pas une contrainte --- c'est une
          \textbf{condition de l'efficacité}~: un dossier propre est
          un dossier solide.

    \item \textbf{Trois cadres éthiques}~: Kant (droits absolus)~;
          Mill (conséquences)~; Aristote (vertus). Les quatre vertus
          de l'investigateur~: intégrité, courage, prudence, justice.

    \item \textbf{Le biais de confirmation est le plus dangereux}~:
          il opère inconsciemment et réinterprète chaque indice
          pour confirmer l'hypothèse initiale. Contre-mesure~:
          E9bis PIU (falsification active).

    \item \textbf{Méthode PEACE}~: Préparation~$\to$ Engage~$\to$
          Account~$\to$ Closure~$\to$ Evaluation. Plus efficace que
          l'interrogatoire coercitif, et juridiquement recevable.

    \item \textbf{Responsabilité personnelle}~: l'investigateur ne
          peut pas se cacher derrière ``j'exécutais un ordre''.
          Violation de procédure~= faute personnelle~; faux~= crime
          personnel~; complicité de violation~= poursuite personnelle.

    \item \textbf{ACPO Principe~4}~: le responsable de l'investigation
          est personnellement responsable de la conformité de toute
          son équipe. Pas de délégation aveugle.

    \item \textbf{Journal de bord éthique}~: documenter les pressions
          reçues, les moments de doute, les décisions difficiles.
          C'est à la fois un outil de développement professionnel
          et une protection légale.

    \item \textbf{Cinq questions de clôture}~: cherché des preuves
          contraires~? Respecté les droits~? Documenté tous les
          transferts~? Noté les pressions~? Rapport défendable
          mot à mot~?
\end{itemize}
\end{boxresume}

\bigskip

\begin{boxexo}[title={\ding{45}\quad Exercice~13.1 --- Identification éthique \niveaubase}]
Pour chacune des situations suivantes, identifiez~: (A)~la tension
déontologique en jeu, (B)~le cadre éthique applicable (Kant / Mill /
Aristote), (C)~l'action correcte~:
\begin{enumerate}[(a)]
    \item Un investigateur découvre, au cours d'une analyse forensique
          d'un PC suspect, des photos compromettantes pour un ministre~---
          sans lien avec l'affaire en cours. Son supérieur lui demande
          de ``signaler informellement'' cette découverte à un tiers.
    \item Un OPJ a passé 3~semaines à construire un dossier contre un
          suspect~; une nouvelle preuve suggère une autre personne.
          Un collègue lui dit~: ``tu as déjà trop de travail dans ce
          dossier pour recommencer''.
    \item L'entreprise victime demande à l'investigateur de ne pas
          inclure dans son rapport les failles de sécurité internes
          qui ont facilité l'attaque~--- pour protéger sa réputation.
\end{enumerate}
\end{boxexo}

\begin{boxexo}[title={\ding{45}\quad Exercice~13.2 --- Biais cognitifs \niveaumoyen}]
Relisez ce fragment de rapport d'investigation~:

\begin{quote}
\textit{``Lors de l'entretien du 12~mars, BELLO Armand a présenté un
comportement nerveux (regards fuyants, mains tremblantes) confirmant
sa culpabilité dans le détournement. L'analyse des logs révèle 3~connexions
depuis son poste, cohérentes avec notre hypothèse. Le fait que d'autres
employés aient eu accès au même système ne constitue pas un obstacle
significatif à notre conclusion.''}
\end{quote}

\begin{enumerate}[(a)]
    \item Identifiez au moins trois biais cognitifs dans ce passage.
    \item Réécrivez ce passage en séparant strictement constatations
          et interprétations, et en appliquant le principe de falsification.
    \item Quelles questions supplémentaires cet investigateur aurait-il
          dû se poser avant de rédiger ce passage~?
\end{enumerate}
\end{boxexo}

\begin{boxexo}[title={\ding{45}\quad Exercice~13.3 --- Simulation décision éthique \niveauexpert}]
Vous êtes l'expert judiciaire désigné dans l'affaire WOURI
(Chapitre~\ref{chap:gr-affaires}). À l'audience, l'avocat de la défense
vous pose les questions suivantes~:
\begin{enumerate}[(a)]
    \item \textit{``Votre rapport affirme que l'email provenait d'un
          serveur russe. Êtes-vous certain à 100\% de cette conclusion~?''}
    \item \textit{``Vous avez utilisé Volatility~3 pour analyser la RAM.
          N'est-il pas vrai que ce logiciel modifie lui-même la RAM
          pendant la capture~?''}
    \item \textit{``Vous êtes mandaté par le parquet. Vos conclusions
          ne sont-elles pas par nature orientées contre mon client~?''}
\end{enumerate}
Rédigez vos réponses telles que vous les donneriez à l'audience, en
respectant les obligations déontologiques de l'expert judiciaire.
Identifiez pour chaque réponse le principe éthique mobilisé.
\end{boxexo}

\bigskip

\begin{boxinfo}
\textbf{Clôture de la Partie~III.} Ce chapitre clôt la Partie~III
(Cadre normatif et déontologie). Les quatre chapitres de cette partie
ont construit votre culture juridique~: cadre normatif global (Ch.~10),
législation mondiale (Ch.~11), droit camerounais (Ch.~12) et déontologie
(ce chapitre).

La \textbf{Partie~IV} aborde les techniques avancées~: forensique
réseau opérationnelle, anti-forensique et contre-mesures, analyse
formelle et cryptanalyse. Ces chapitres mobilisent l'intégralité
des fondements posés dans les Parties~I à~III~--- techniques,
juridiques et déontologiques.
\end{boxinfo}
